\chapter{Desarrollo e implementación}
\label{ch:desarrollo_implementacion}

Este capítulo describe cómo se ha materializado la solución propuesta en una
base de código ejecutable. Mientras que el capítulo anterior se centraba en el
diseño de la arquitectura, el motor de recomendación y el modelo de datos, aquí
el foco se desplaza hacia la implementación efectiva: entorno de trabajo,
tecnologías concretas, estructura del código, decisiones de integración ya
aterrizadas en servicios reales y proceso de puesta en marcha. Por ello, este
capítulo evita repetir la justificación conceptual del diseño y se centra en
cómo quedó plasmado en el repositorio final.

\section{Entorno de desarrollo}
\label{sec:entorno_desarrollo}

El desarrollo se ha realizado en un entorno local basado en macOS sobre
arquitectura ARM, con terminal \texttt{zsh}, control de versiones con Git y una
separación clara entre frontend móvil y backend API. La aplicación móvil se ha
trabajado sobre Flutter estable, mientras que el backend se ha mantenido en un
entorno virtual de Python independiente para aislar dependencias.

\begin{table}[H]
\centering
\scriptsize
\renewcommand{\arraystretch}{1.1}
\setlength{\tabcolsep}{4pt}
\begin{tabularx}{\textwidth}{|P{3.0cm}|P{3.0cm}|Y|}
\hline
\textbf{Elemento} & \textbf{Versión o entorno} & \textbf{Papel en el proyecto} \\
\hline
Sistema base & macOS / Darwin 25.3.0 arm64 & Desarrollo, pruebas locales y compilación de memoria y app. \\
\hline
Flutter SDK & Flutter 3.41.2 & Desarrollo del cliente móvil multiplataforma. \\
\hline
Dart SDK & Dart 3.11.0 & Lenguaje principal del frontend. \\
\hline
Python & Python 3.13.1 & Lenguaje principal del backend y del entrenamiento del modelo. \\
\hline
Servidor ASGI & Uvicorn 0.34.0 & Ejecución local del backend FastAPI. \\
\hline
Control de versiones & Git 2.45.1 & Seguimiento de cambios y evolución del código fuente. \\
\hline
Entorno Python & \texttt{.venv} local & Aislamiento de dependencias del backend y del pipeline ML. \\
\hline
\end{tabularx}
\caption{Entorno de desarrollo utilizado durante la implementación.}
\label{tab:entorno_desarrollo}
\end{table}

La elección de este entorno responde a una necesidad práctica: poder iterar con
rapidez sobre frontend, backend y modelo supervisado sin mezclar dependencias
ni comprometer la reproducibilidad del proyecto.

\section{Tecnologías utilizadas}
\label{sec:tecnologias}

Harmony Hub combina un cliente móvil declarativo, un backend de servicios y una
capa documental en la nube. La selección tecnológica no se hizo buscando el
mayor número posible de herramientas, sino un conjunto razonablemente compacto
que cubriera movilidad, persistencia, integración externa y aprendizaje
supervisado.

\begin{table}[H]
\centering
\scriptsize
\renewcommand{\arraystretch}{1.08}
\setlength{\tabcolsep}{2.8pt}
\begin{tabularx}{\textwidth}{|P{2.45cm}|P{3.2cm}|Y|}
\hline
\textbf{Tecnología} & \textbf{Uso en el proyecto} & \textbf{Justificación} \\
\hline
Flutter + Dart & Interfaz móvil, navegación, gestión de estado local y acceso a servicios de cliente. & Permite desarrollar una app móvil completa con una sola base de código y una iteración visual muy rápida. \\
\hline
FastAPI & API backend, orquestación del motor de recomendación e integración con Spotify. & Ofrece una estructura clara de rutas y servicios, baja fricción de desarrollo y buen encaje con Python. \\
\hline
Cloud Firestore & Persistencia documental de usuarios, check-ins, recomendaciones, playlists y feedback. & Encaja bien con el modelo orientado a documentos del proyecto y simplifica la sincronización entre cliente y backend. \\
\hline
Firebase Authentication & Gestión de identidad del usuario. & Reduce complejidad en autenticación básica y aporta integración directa con el ecosistema Firebase. \\
\hline
Spotify Web API & Obtención de perfil musical básico, autorización PKCE y creación de playlists reales. & Es la pieza que convierte la recomendación en una salida utilizable dentro del ecosistema musical cotidiano del usuario sin asumir la fase principal de ranking. \\
\hline
scikit-learn & Entrenamiento y validación del modelo supervisado de sesión. & Proporciona herramientas maduras para preprocesado, validación cruzada, métricas y clasificación probabilística. \\
\hline
\shortstack[l]{Cryptography +\\Flutter Secure Storage} & Cifrado del bloque emocional y almacenamiento local de claves. & Reducen la exposición de datos emocionales sensibles y ayudan a separar parte pública y privada de la sesión. \\
\hline
\shortstack[l]{Noise Meter +\\Permission Handler} & Medición opcional del entorno acústico en el dispositivo. & Permiten capturar contexto ambiental sin grabar ni persistir audio bruto. \\
\hline
\shortstack[l]{Dio +\\HTTPX} & Comunicación HTTP entre cliente y backend. & Mantienen una capa de comunicación explícita y fácil de depurar durante el desarrollo. \\
\hline
\end{tabularx}
\caption{Tecnologías principales empleadas en la implementación de Harmony Hub.}
\label{tab:tecnologias}
\end{table}

\section{Implementación de los módulos principales}
\label{sec:implementacion_modulos}

La implementación sigue la separación funcional introducida en el diseño, pero
traducida a una estructura de carpetas y responsabilidades concretas.

\subsection{Cliente móvil Flutter}
\label{subsec:modulo_cliente_flutter}

El cliente se organiza por features. Dentro de \texttt{harmonyhub/lib},
cada bloque funcional dispone de sus propias capas de \texttt{data} y
\texttt{presentation}: autenticación, check-in, recomendación, Spotify,
feedback, historial, modos predefinidos y aprendizaje del usuario. Esta
organización favorece que cada parte de la app pueda evolucionar con cierto
aislamiento sin depender de una carpeta única de pantallas o servicios.

En la práctica, el frontend implementa:

\begin{itemize}
    \item el flujo de acceso y validación de sesión;
    \item la captura guiada del check-in con preguntas y respuestas cotidianas;
    \item la medición opcional del entorno acústico desde el dispositivo;
    \item la visualización de la recomendación previa a la generación real;
    \item la apertura del flujo de creación de playlist y recogida de feedback;
    \item la lectura de historial, modos y aprendizaje acumulado.
\end{itemize}

Una decisión de implementación importante ha sido mantener la lógica de red y
persistencia fuera de las pantallas. Para ello, las vistas se apoyan en
servicios específicos de check-in, recomendación, feedback y Spotify, de forma
que la interfaz no concentre también la responsabilidad de hablar con backend o
con Firestore.

\subsection{Backend API y orquestación}
\label{subsec:modulo_backend_api}

El backend se articula en dos niveles. El primero es una capa de rutas HTTP
definida en \texttt{backend\_fastapi/app/api/routes}, donde cada recurso
principal cuenta con su propio punto de entrada:

\begin{itemize}
    \item \texttt{/checkins} para registrar sesiones iniciales;
    \item \texttt{/recommendations} para generar la recomendación de sesión;
    \item \texttt{/spotify} para autenticación, perfil y creación de playlists;
    \item \texttt{/feedback} para cerrar la sesión y alimentar aprendizaje;
    \item \texttt{/history} para exponer, cuando se necesita desde backend, una
    vista del historial persistido real del usuario.
\end{itemize}

El segundo nivel es una capa de servicios ubicada en el directorio de servicios
del backend. Aquí reside la mayor parte del comportamiento del sistema: motor
heurístico de recomendación, generación de playlists, integración semántica,
similitud vectorial, aprendizaje del usuario, extracción de ejemplos ML y
comunicación con Firestore.

Esta estructura evita que la ruta HTTP se convierta en un bloque monolítico. La
ruta recibe, valida y delega; el servicio interpreta y ejecuta; la persistencia
almacena o recupera. Ese reparto facilita depuración, ampliación del sistema y
reutilización interna.

\subsection{Motor de recomendación y generación de playlist}
\label{subsec:modulo_recomendacion_implementacion}

En la implementación real, el motor de recomendación no se reduce a una única
función. Intervienen varios servicios con papeles distintos:

\begin{itemize}
\item \texttt{recommender\_service.py}, centrado en la generación y
    puntuación de candidatos de modo de sesión y en la abstención controlada
    del modelo supervisado cuando no supera sus puertas de datos o calidad, o
    cuando la sesión concreta no alcanza confianza suficiente;
    \item el servicio de modelo profesional de playlist, responsable de
    transformar el modo elegido en un perfil musical operativo;
    \item los servicios de generación personalizada, análisis semántico y
    similitud vectorial, que ayudan a recuperar, valorar y reordenar canciones
    candidatas.
\end{itemize}

En términos de código, la recomendación pasa por tres pasos encadenados:
selección del modo de sesión, construcción del generation profile y
materialización de la playlist. Entre el segundo y el tercer paso se apoya
además en el catálogo msd\_tracks, que en el estado actual reúne 5\,441
pistas y funciona como fuente principal de ranking. Esa división ha permitido
modificar reglas, catálogo, aprendizaje individual o integración con Spotify
sin rehacer el pipeline completo en cada iteración.

\subsection*{Lectura de una traza típica de generación}

Una ventaja práctica de la implementación es que el backend deja trazas lo
bastante expresivas como para reconstruir el recorrido completo de una sesión
musical. Esto resulta útil tanto para depuración como para validación, ya que
permite observar cómo el sistema pasa de los candidatos de modo a la playlist
realmente creada. Para ello se toma como referencia una ejecución completa
realizada desde un dispositivo Android físico conectado a la misma red local
que el backend, con autenticación real en Spotify, medición de entorno activa y
cierre correcto del flujo con feedback.

Para concretar el comportamiento descrito, resulta útil fijar un ejemplo
completo de extremo a extremo para mostrar cómo se encadenan realmente las
distintas capas del sistema. En la validación final se registró una sesión de
relajación con usuario en estado neutral, estrés medio, energía media,
preferencia instrumental, intensidad suave, exploración orientada a descubrir,
popularidad alternativa y medición de entorno activa. El entorno inferido fue
\texttt{Entorno mixto}, con categoría de ruido \texttt{loud}, confianza
\texttt{0.95}, variabilidad \texttt{2.7474}, diferencia de picos
\texttt{13.2875}, proporción transitoria \texttt{0.4333} y
\texttt{burst\_count = 4}. Sobre esa base, el backend dejó una traza completa
que permite reconstruir el flujo general en orden.

\subsection*{Ejemplo ordenado de flujo completo}

En una primera fase, el sistema resuelve la autenticación con Spotify mediante
OAuth con PKCE. Las líneas \texttt{[AUTH] login\_url\_created},
\texttt{authorize\_url: ...} y \texttt{state: ...} indican que el backend ha
generado una URL de autorización con \texttt{client\_id},
\texttt{redirect\_uri}, \texttt{scope}, \texttt{state} y
\texttt{code\_challenge}. Después, el bloque
\texttt{[AUTH] exchange\_received} seguido de
\texttt{[AUTH] exchange\_code\_for\_token} muestra que la aplicación móvil ya
ha recibido el \textit{authorization code} y lo está intercambiando por un
\textit{access token}. Finalmente, \texttt{[AUTH] token\_received} y
\texttt{[AUTH] token\_obtained} confirman que el backend dispone de un token
\texttt{Bearer} válido, con alcance suficiente para consultar el perfil del
usuario y crear playlists reales en su cuenta.

En segundo lugar, la aplicación valida ese token y recupera el perfil del
usuario. Las trazas \texttt{[PROFILE] me\_received},
\texttt{[PROFILE] get\_spotify\_profile} y
\texttt{[PROFILE] profile\_obtained} confirman que el backend ha consultado el
endpoint \texttt{/me} de Spotify y ha recuperado correctamente el
\texttt{spotify\_user\_id}, el \texttt{display\_name} y el mercado del usuario.
Este paso es importante porque vincula la sesión actual con una identidad
musical real, sobre la que después se materializará la playlist.

En tercer lugar, una vez almacenado el check-in (\texttt{POST /checkins/}),
entra en funcionamiento el clasificador supervisado de modo de sesión. En esta
ejecución, el backend indica \texttt{[ML] model loaded} y
\texttt{[ML] candidates: 3}, lo que significa que el modelo compara tres
modos plausibles de la misma familia funcional:
\texttt{relajacion\_neutral\_suave},
\texttt{relajacion\_neutral\_media} y
\texttt{relajacion\_neutral\_alta}. El diccionario mostrado en
\texttt{[ML] context: ...} es especialmente revelador porque deja visible el
conjunto completo de señales que alimentan al modelo: objetivo, mood, niveles
de estrés y energía, preferencia vocal, intensidad, exploración, popularidad,
duración deseada y variables ambientales derivadas del micrófono.

Las probabilidades \texttt{[0.9454, 0.9479, 0.9479]} muestran que el modelo
considera muy viables los tres candidatos. El valor
\texttt{probability\_spread = 0.0025} se calcula como la diferencia entre la
probabilidad máxima y la mínima, por lo que aquí expresa una separación muy
pequeña entre alternativas. Esto significa que el ML está activo y confía en
los candidatos, aunque no encuentra una diferencia supervisada fuerte entre
ellos. El hecho de que \texttt{selected\_mode\_probability = 0.9479} supere
con amplitud el umbral \texttt{min\_selected\_mode\_probability = 0.14}
explica que el sistema no entre en abstención y permita a la capa supervisada
ajustar la decisión heurística.

En cuarto lugar aparece el bloque \texttt{[SESSION]}, donde ya no se observan
probabilidades aisladas sino la decisión híbrida final. En esta traza,
\texttt{ml\_enabled = True} confirma que la capa supervisada ha pasado las
puertas globales necesarias. Cada línea \texttt{CANDIDATE} combina tres
elementos: una puntuación heurística base, una probabilidad del modelo y un
\texttt{ml\_delta} que transforma esa probabilidad en ajuste numérico. En el
caso \texttt{relajacion\_neutral\_suave}, la heurística inicial ya era muy
superior (\texttt{46.0}) a la de los modos vecinos (\texttt{17.0} y
\texttt{6.0}); el ML añade a ese candidato un incremento de \texttt{8.91} y
eleva la puntuación final a \texttt{54.91}. Aunque los modos medio y alto
obtienen probabilidades prácticamente idénticas, su base heurística era mucho
menor, de modo que el candidato suave conserva la primera posición. Por eso la
traza final indica \texttt{selected\_mode = relajacion\_neutral\_suave} y
\texttt{selection\_source = session\_ml}: el ML participa de forma real, pero
lo hace reforzando una ventaja heurística ya bien fundamentada.

En quinto lugar, el bloque \texttt{[ENV]} permite comprobar que el entorno no
actúa como una variable decorativa, sino como parte operativa de la
personalización. El sistema marca \texttt{use\_environment = True},
\texttt{environment\_context = Entorno mixto} y
\texttt{environment\_influence\_strength = 0.95}. Esta última magnitud se
deduce a partir de la confianza, la variabilidad, la diferencia de picos, la
proporción transitoria y el número de ráfagas detectadas, y en este caso queda
saturada en el máximo operativo previsto por el sistema. El aspecto más
significativo de la traza es el campo \texttt{target\_adjustment}: el perfil
musical previo partía de \texttt{valence = 0.55}, \texttt{energy = 0.12} y
\texttt{danceability = 0.10}, pero tras aplicar el ruido \texttt{loud} dentro
de un objetivo de relajación, la energía pasa a \texttt{0.348}. Es decir, el
sistema entiende que una relajación eficaz en un ambiente sonoro fuerte no debe
basarse en música excesivamente tenue, sino en una calma con algo más de
presencia. A esto se añade el bloque \texttt{noise\_queries}, que incorpora
consultas como \texttt{stable calm focus}, \texttt{warm ambient} o
\texttt{soft but present chill}. En cambio, \texttt{context\_queries} aparece
vacío porque el contexto \texttt{Entorno mixto} no activa una familia
específica de consultas contextuales adicionales.

En sexto lugar, el bloque \texttt{[LEARNING]} documenta el estado del
aprendizaje musical acumulado para el mood actual. La línea
\texttt{mood\_gate\_passed = True} indica que el sistema ya dispone de
evidencia suficiente para aplicar aprendizaje específico en el mood neutral. El
valor \texttt{mood\_quality\_score = 81.63} resume consistencia, fuerza
observacional y cobertura de géneros positivos, mientras que
\texttt{mood\_application\_factor = 1.0} expresa que ese aprendizaje no se
aplica de forma parcial, sino completa.

En séptimo lugar, tras cerrarse correctamente la recomendación
(\texttt{POST /recommendations/generate ... 200 OK}), el backend entra en la
fase de generación musical propiamente dicha. El bloque \texttt{[INPUT]} fija
de nuevo el contexto de trabajo mediante \texttt{recommendation\_id},
\texttt{goal}, \texttt{mood}, \texttt{stress}, \texttt{energy},
\texttt{outcome}, uso del entorno y preferencias declaradas. A partir de ahí,
la línea \texttt{[LEARNING] mode=... | subtype=... | curve=... |
profile\_mode=...} sintetiza el \textit{generation profile}. En el ejemplo
validado, el subtipo resultante es \texttt{stable\_relaxation}, la curva de
activación es \texttt{peak\_then\_settle} y el modo de perfil es
\texttt{stable\_weighted}, con \texttt{session\_weight = 0.169} y
\texttt{stable\_weight = 0.225}. Esto significa que el sistema sigue teniendo
en cuenta la señal reciente de la sesión, pero se apoya algo más en el patrón
estable aprendido del usuario. En la implementación, esta capa de
subtipos se ha ampliado además con variantes específicas de voz, como
\texttt{guided\_focus} para sesiones de foco con presencia vocal controlada y
\texttt{comfort\_relaxation} para sesiones de relajación con voz cálida y tono
de acompañamiento. De este modo, la preferencia vocal ya no opera solo como
filtro de canciones, sino también como parte estructural del perfil musical
interno.

En octavo lugar, la traza entra en la fase de selección de canciones. La línea
\texttt{[AFFINITY] source = msd\_only} recuerda que la
fuente principal del ranking es el catálogo estructurado local y no una
recuperación dinámica viva desde Spotify. A continuación,
\texttt{[CATALOG] dataset\_candidates = 60 | top\_track = Sunset Lover - Petit
Biscuit | top\_score = 123.85} resume el estado del ranking principal: sesenta
canciones han entrado como candidatas desde el catálogo MSD, y la pista mejor
puntuada resulta coherente con una sesión de relajación instrumental y
alternativa.

En noveno lugar se produce la materialización sobre Spotify. El bloque
\texttt{[MATERIALIZATION] playable = 15 | cached = 1 | fresh = 14 |
failed = 2 | searches = 16 | rate\_limited = False} indica que quince
candidatos del ranking pudieron vincularse a URIs reales de Spotify, uno ya
estaba cacheado, catorce se resolvieron mediante búsquedas nuevas y dos no
llegaron a materializarse. Esta capa es importante porque separa claramente dos
problemas distintos: el sistema puede rankear bien dentro del catálogo local y,
después, todavía tener que convertir ese ranking en canciones realmente
reproducibles en la plataforma externa.

En décimo lugar aparece el ensamblado final de la sesión:
\texttt{[PLAYLIST] selected\_tracks = 10 | uris = 10 | target\_ms = 2400000 |
actual\_ms = 2725096 | missing\_duration = 0}. Esto significa que la playlist
se ha construido con diez canciones válidas, todas con URI disponible, y con
una duración final aproximada de 45 minutos frente a un objetivo inicial de 40.
Aunque la duración final supera el objetivo, el sistema no deja huecos sin
resolver ni canciones sin duración conocida. Después, el bloque
\texttt{[MATERIALIZATION] add\_items\_to\_playlist} seguido de
\texttt{playlist\_id = ...} y \texttt{uris\_count = 10} marca la escritura
real en Spotify, y la línea \texttt{[PLAYLIST] created id = ... |
name = Harmony Hub · Relajacion | query\_signal\_count = 2} cierra la sesión
confirmando que la salida final ya no es una recomendación abstracta, sino una
playlist real creada en la cuenta autenticada del usuario.

Por último, el flujo se cierra con feedback y mantenimiento automático del
modelo. La respuesta \texttt{POST /feedback/ ... 200 OK} confirma que la
valoración de la sesión se ha persistido correctamente. Inmediatamente después,
el backend deja las trazas \texttt{[ML MAINTENANCE] start},
\texttt{train} y \texttt{done}, indicando que se ha detectado un nuevo ejemplo
supervisado (\texttt{labeled\_examples = 65}) y que el sistema ha ejecutado el
reentrenamiento posterior. Así, una única sesión queda enlazada con todo el
ciclo del producto: autenticación, check-in, recomendación híbrida, uso
efectivo del entorno, ranking sobre catálogo local, materialización en
Spotify, feedback del usuario y actualización posterior del componente
supervisado.

Esta lectura deja ver con bastante claridad la separación entre capas del
sistema. Primero se decide el modo de sesión, después se construye el perfil
musical operativo, a continuación se consulta y reordena el catálogo local, y
solo al final se materializa la playlist en Spotify. El cierre con feedback
activa además una segunda traza, ya no de recomendación musical, sino de
mantenimiento del aprendizaje supervisado, lo que permite enlazar una misma
sesión con su efecto posterior sobre el perfil y sobre el modelo.

\subsection{Persistencia, historial y aprendizaje}
\label{subsec:modulo_persistencia_aprendizaje}

La persistencia combina escrituras directas desde el cliente y escrituras
derivadas desde backend. El cliente registra check-ins, recomendaciones
visibles, playlists generadas y feedback, mientras que el backend consolida
estructuras agregadas para aprendizaje, trazabilidad y entrenamiento.

En términos de implementación, esto permite dos ventajas:

\begin{itemize}
    \item mantener disponible el historial funcional de la app aunque una
    sesión concreta termine en abstención local del modelo;
    \item separar la experiencia operativa del usuario de los artefactos
    internos de entrenamiento y aprendizaje.
\end{itemize}

El historial reconstruye sesiones enlazando documentos públicos y privados,
mientras que el aprendizaje progresivo deriva patrones de uso y ejemplos de
entrenamiento a partir del ciclo cerrado de recomendación, escucha y feedback.
En la versión final se corrigió además una inconsistencia arquitectónica
importante: el historial visible de la app ya no depende de colecciones
temporales en memoria del backend, sino de documentos persistidos en Firestore.
La interfaz Flutter reconstruye el histórico directamente desde esas
colecciones, y el endpoint \texttt{/history/me} quedó alineado para exponer esa
misma fuente persistida en su parte pública.

\section{Integraciones y servicios externos}
\label{sec:integraciones}

Las integraciones externas ya justificadas en el diseño se materializan en tres
frentes concretos:

\begin{itemize}
    \item \textbf{Firebase Authentication}, usado como origen de identidad
    estable para el cliente y como base sobre la que se apoyan las reglas de
    acceso en Firestore.
    \item \textbf{Cloud Firestore}, empleado como persistencia documental
    operativa tanto para el flujo visible de usuario como para estructuras de
    aprendizaje, ejemplos supervisados y trazabilidad.
    \item \textbf{Spotify Web API}, utilizada para OAuth con PKCE, perfil
    musical básico del usuario y materialización final de playlists.
\end{itemize}

La implementación intenta además que estas dependencias no se conviertan
en puntos únicos de fallo. El backend sigue pudiendo rankear sobre el catálogo
MSD aunque Spotify limite búsquedas o retrase la materialización, y el cliente conserva historial y
documentos propios incluso cuando una sesión concreta termina en abstención local
del clasificador por falta de confianza.

\subsection*{Medición del entorno acústico}

La app utiliza el micrófono del dispositivo únicamente para obtener una ventana
breve de medición y derivar un perfil agregado del entorno. La implementación no
persiste audio en bruto ni fragmentos grabados; solo conserva métricas
resumidas, como media, mediana, rango, variabilidad o contexto inferido.
Además, se añadieron dos señales de calidad de medición:
\texttt{sample\_density\_hz} y \texttt{stability\_score}. A partir de ellas, y
de una \texttt{confidence} agregada, el sistema calcula una intensidad de
influencia ambiental que no se limita a activar o desactivar el entorno, sino
que pondera cuánto debe pesar en la sesión. Esa decisión se persiste después
mediante campos como \texttt{environment\_usage\_status} y
\texttt{environment\_usage\_rationale}, lo que mejora bastante la trazabilidad
del sistema.

La evidencia operativa recogida durante la validación con dispositivo Android
físico muestra sesiones con \texttt{use\_environment = True}, contexto
ambiental inferido, fuerza de influencia positiva y consultas auxiliares
derivadas del entorno medido. Esto permite afirmar que la medición acústica no
se limita a una etiqueta visible en interfaz o a un dato almacenado en
Firestore, sino que participa realmente en la construcción del perfil musical,
en el refuerzo de consultas y en la selección final de canciones de la
playlist.

\subsection*{Pipeline de entrenamiento supervisado}

El modelo de sesión se entrena fuera del flujo interactivo principal, pero ya
no depende exclusivamente de que un desarrollador ejecute
comandos manuales. Tras cada feedback útil para entrenamiento, el backend
dispara una tarea de mantenimiento que audita el número de ejemplos etiquetados,
compara el estado actual con los metadatos previos y decide si procede
reentrenar. Cuando se detectan ejemplos nuevos, o incluso un reinicio o
reducción del dataset respecto a metadatos antiguos, el sistema vuelve a
entrenar automáticamente y actualiza dos artefactos: el modelo serializado y un
fichero de auditoría.

El script \texttt{train\_ranking\_model.py} sigue existiendo como mecanismo
manual de recuperación o depuración, pero el flujo ordinario del prototipo ya
queda cubierto por ese mantenimiento automático. En ambos casos, la fase de
entrenamiento calcula observabilidad por mood, estado de puertas de calidad y
motivo de disponibilidad o abstención del modelo.

En el estado actual del repositorio, este pipeline ya genera no solo el modelo
serializado y sus metadatos, sino también artefactos de auditoría y
explicabilidad. La situación real del clasificador global ya no es de
abstención completa: dispone de 65 ejemplos de sesión, con un balance
suficiente para la configuración operativa final, puerta global activa,
cobertura completa por mood y umbral de confianza configurable. Esta
observación es importante porque refleja que el pipeline no solo está terminado
e integrado, sino también activado en ejecución real dentro del flujo del
producto.

\section{Seguridad, privacidad y tratamiento de datos}
\label{sec:seguridad_privacidad}

La implementación incorpora varias medidas orientadas a reducir la exposición de
datos personales y emocionales.

\subsection*{Separación entre bloque público y bloque privado}

Los campos más sensibles del check-in y del feedback no se guardan únicamente en
el documento funcional principal. Se separan en colecciones privadas paralelas
(\texttt{checkins\_private} y \texttt{feedback\_private}), mientras que el
documento público conserva solo el mínimo necesario para trazabilidad y uso de
la aplicación.

\subsection*{Cifrado en cliente del bloque emocional}

El cliente móvil cifra la carga emocional usando AES-GCM de 256 bits antes de
persistirla. La clave asociada al usuario se almacena localmente mediante
\texttt{flutter\_secure\_storage}, y el bloque cifrado conserva
\texttt{ciphertext}, \texttt{nonce} y \texttt{mac}. Esta decisión no sustituye
por completo otras capas de seguridad, pero sí reduce exposición innecesaria de
estados emocionales explícitos.

\subsection*{Control de acceso en Firestore}

Las reglas de Firestore restringen el acceso por propietario y validan la
forma de los documentos aceptados. Esto se aprecia especialmente en las
colecciones privadas y en los documentos que referencian al usuario autenticado.
De este modo, el proyecto no delega toda la protección en el cliente, sino que
la refuerza también en la política de acceso documental.

En concreto, \texttt{firestore.rules} combina comprobaciones de identidad y de
esquema. Funciones como \texttt{isOwnerByField()},
\texttt{isRequestOwnerByField()} e \texttt{isOwnerByDocId(docId)} impiden que
un usuario lea o escriba documentos ajenos, mientras que
\texttt{onlyHasKeys(...)} fuerza que colecciones como \texttt{checkins},
\texttt{feedback} y \texttt{generated\_playlists} solo acepten los campos
previstos por el diseño. Además, las colecciones
\texttt{user\_generation\_preferences}, \texttt{training\_examples} y
\texttt{training\_session\_examples} permanecen cerradas a escritura desde el
cliente, de modo que el aprendizaje agregado y los ejemplos supervisados quedan
reservados al backend o a cuentas con privilegios administrativos.

\subsection*{Gestión de permisos y minimización de datos}

El acceso al micrófono se solicita únicamente cuando el usuario activa la
medición del entorno, y la salida persistida se limita a métricas resumidas del
ambiente. De forma análoga, la integración con Spotify se activa solo tras la
autorización explícita del usuario y dentro del alcance que necesita la app
para leer contexto musical y crear playlists.

En conjunto, estas decisiones no eliminan por completo el riesgo asociado al
tratamiento de información sensible, pero sí reducen de forma apreciable la
superficie de exposición. La separación entre bloque funcional y bloque privado,
el cifrado previo a la persistencia, las reglas de acceso por propietario y la
minimización del dato capturado permiten que Harmony Hub trate estados
emocionales y valoraciones de sesión con un criterio de prudencia técnica
coherente con la naturaleza del problema abordado.

\section{Síntesis operativa del proyecto}
\label{sec:estado_actual_proyecto}

En el momento de cierre de esta memoria, Harmony Hub puede describirse ya como
una solución funcional completa y trazable de extremo a extremo. Esta
conclusión se apoya en tres hechos verificados durante la fase final de
validación.

En primer lugar, los datos de entrada del usuario están controlados y
formalizados. El contexto de sesión ya no depende de descripciones libres o de
interpretaciones implícitas, sino de un check-in guiado con opciones cerradas
para estado emocional, objetivo, energía, estrés, preferencias musicales,
duración y resultado esperado. A esta capa se añade, cuando el usuario activa
la medición, un bloque acústico estructurado con categoría de ruido, contexto
ambiental, confianza, estabilidad y otras señales agregadas. El resultado es un
contexto consistente que puede validarse, persistirse y reutilizarse en todo el
pipeline.

En segundo lugar, el entorno acústico ya puede defenderse como parte efectiva
de la personalización. La validación ha mostrado sesiones reales con uso del
entorno activado, contexto ambiental inferido, fuerza de influencia elevada y
consultas auxiliares no vacías asociadas tanto al ruido como al contexto
detectado. Por tanto, el entorno no se queda en una función accesoria de
interfaz, sino que influye de forma operativa en la recomendación y en la
generación final de la playlist.

En tercer lugar, el flujo completo resulta trazable desde el check-in hasta el
feedback. La persistencia en Firestore enlaza documentos de sesión,
recomendación, playlist y valoración posterior, mientras que el backend añade
ejemplos supervisados de sesión, metadatos del clasificador y auditoría de
mantenimiento automático. Esta doble traza, funcional y técnica, permite
explicar no solo qué playlist se generó, sino también con qué contexto, con qué
modo, con qué peso del entorno y con qué consecuencias posteriores sobre el
aprendizaje del sistema.

Desde la perspectiva de entrega, el estado actual del repositorio permite
sostener que Harmony Hub ya no es una prueba de concepto desconectada, sino una
arquitectura híbrida integrada donde conviven captura de contexto,
recomendación explicable, uso real del entorno, activación supervisada cuando
existe evidencia suficiente, materialización completa en Spotify y cierre
posterior del ciclo de aprendizaje.

\section{Despliegue o puesta en marcha}
\label{sec:despliegue}

La solución está preparada para ejecutarse en local separando frontend, backend
y reglas documentales.

\subsection*{Puesta en marcha del frontend}

En el entorno de desarrollo, el cliente móvil puede arrancarse con los
comandos:

\begin{minted}[fontsize=\small, breaklines]{bash}
cd harmonyhub
flutter pub get
flutter run
\end{minted}

Este proceso instala dependencias, resuelve paquetes y ejecuta la aplicación en
emulador o dispositivo conectado.

\subsection*{Puesta en marcha del backend}

El backend se levanta en un entorno virtual independiente:

\begin{minted}[fontsize=\small, breaklines]{bash}
cd backend_fastapi
python -m venv .venv
source .venv/bin/activate
pip install -r requirements.txt
uvicorn app.main:app --reload
\end{minted}

Con ello se expone la API REST necesaria para check-ins, recomendaciones,
feedback, historial y flujo Spotify. En la validación entregable, la
aplicación se utilizó instalada en un dispositivo Android físico y conectada a
un backend activo en la misma red local.

\subsection*{Entrenamiento del modelo supervisado}

En el uso normal del sistema, el modelo ya se audita y se reentrena
automáticamente en segundo plano tras el cierre del feedback, siempre que haya
ejemplos nuevos o que el dataset actual haya quedado desalineado respecto a los
metadatos previos. Aun así, para tareas de mantenimiento o depuración sigue
siendo posible lanzar el entrenamiento manualmente mediante:

\begin{minted}[fontsize=\small, breaklines]{bash}
cd backend_fastapi
source .venv/bin/activate
python -m app.ml.train_ranking_model
\end{minted}

Este proceso genera el artefacto serializado del modelo y un fichero de
metadatos que informa de métricas, volumen de datos, observabilidad por
mood y estado de las puertas de activación. El mantenimiento automático
deja además un segundo artefacto de auditoría que registra si hubo
reentrenamiento, por qué motivo se ejecutó o por qué se omitió.

En consecuencia, el entrenamiento manual queda reservado a recuperación,
depuración o demostración académica del pipeline, mientras que el uso ordinario
de la app mantiene el modelo alineado con el conjunto real de ejemplos sin
intervención explícita del usuario final.

\subsection*{Reglas de Firestore}

La política documental del proyecto también forma parte de la puesta en marcha.
Las reglas se asocian desde \texttt{firebase.json} al archivo
\texttt{firestore.rules} y pueden desplegarse con:

\begin{minted}[fontsize=\small, breaklines]{bash}
firebase deploy --only firestore:rules
\end{minted}
